\documentclass[10pt,a4paper]{amsart}
\usepackage{amsmath,amssymb,amsthm,graphicx}
\usepackage{verbatim}
\newenvironment{reference}{\comment}{\endcomment}
%\newenvironment{reference}{}{}
\newenvironment{slogan}{\comment}{\endcomment}
\newenvironment{history}{\comment}{\endcomment}

% For commutative diagrams we use Xy-pic
\usepackage[all]{xy}

% We use 2cell for 2-commutative diagrams.
\xyoption{2cell}
% Theorem environments.
%
\theoremstyle{plain}
\newtheorem{theorem}[subsection]{Theorem}
\newtheorem{proposition}[subsection]{Proposition}
\newtheorem{lemma}[subsection]{Lemma}

\theoremstyle{definition}
\newtheorem{definition}[subsection]{Definition}
\newtheorem{example}[subsection]{Example}
\newtheorem{exercise}[subsection]{Exercise}
\newtheorem{situation}[subsection]{Situation}

\theoremstyle{remark}
\newtheorem{remark}[subsection]{Remark}
\newtheorem{remarks}[subsection]{Remarks}

\numberwithin{equation}{subsection}

% Macros
%
\def\lim{\mathop{\mathrm{lim}}\nolimits}
\def\colim{\mathop{\mathrm{colim}}\nolimits}
\def\Spec{\mathop{\mathrm{Spec}}}
\def\Hom{\mathop{\mathrm{Hom}}\nolimits}
\def\Ext{\mathop{\mathrm{Ext}}\nolimits}
\def\SheafHom{\mathop{\mathcal{H}\!\mathit{om}}\nolimits}
\def\SheafExt{\mathop{\mathcal{E}\!\mathit{xt}}\nolimits}
\def\Sch{\mathit{Sch}}
\def\Mor{\mathop{\mathrm{Mor}}\nolimits}
\def\Ob{\mathop{\mathrm{Ob}}\nolimits}
\def\Sh{\mathop{\mathit{Sh}}\nolimits}
\def\NL{\mathop{N\!L}\nolimits}
\def\CH{\mathop{\mathrm{CH}}\nolimits}
\def\proetale{{pro\text{-}\acute{e}tale}}
\def\etale{{\acute{e}tale}}
\def\QCoh{\mathit{QCoh}}
\def\Ker{\mathop{\mathrm{Ker}}}
\def\Im{\mathop{\mathrm{Im}}}
\def\Coker{\mathop{\mathrm{Coker}}}
\def\Coim{\mathop{\mathrm{Coim}}}

% Boxtimes
%
\DeclareMathSymbol{\boxtimes}{\mathbin}{AMSa}{"02}

%
% Macros for moduli stacks/spaces
%
\def\QCohstack{\mathcal{QC}\!\mathit{oh}}
\def\Cohstack{\mathcal{C}\!\mathit{oh}}
\def\Spacesstack{\mathcal{S}\!\mathit{paces}}
\def\Quotfunctor{\mathrm{Quot}}
\def\Hilbfunctor{\mathrm{Hilb}}
\def\Curvesstack{\mathcal{C}\!\mathit{urves}}
\def\Polarizedstack{\mathcal{P}\!\mathit{olarized}}
\def\Complexesstack{\mathcal{C}\!\mathit{omplexes}}
% \Pic is the operator that assigns to X its picard group, usage \Pic(X)
% \Picardstack_{X/B} denotes the Picard stack of X over B
% \Picardfunctor_{X/B} denotes the Picard functor of X over B
\def\Pic{\mathop{\mathrm{Pic}}\nolimits}
\def\Picardstack{\mathcal{P}\!\mathit{ic}}
\def\Picardfunctor{\mathrm{Pic}}
\def\Deformationcategory{\mathcal{D}\!\mathit{ef}}
\usepackage{cleveref}
\crefname{originallemma}{original lemma}{original lemmas}
\Crefname{originallemma}{Original lemma}{Original lemmas}
\makeatletter
\newlabel{chow-lemma-chow-ring-projective-bundle}{{02TX}{source}{}{sourcepointer-1}{}}
\newlabel{chow-lemma-chow-ring-projective-bundle@cref}{{[originallemma][0][]02TX}{[1][0][]source}{}{}{}}
\makeatother
\crefname{originallemma}{original lemma}{original lemmas}
\Crefname{originallemma}{Original lemma}{Original lemmas}
\makeatletter
\newlabel{chow-lemma-lci-gysin-flat}{{0FF4}{source}{}{sourcepointer-2}{}}
\newlabel{chow-lemma-lci-gysin-flat@cref}{{[originallemma][0][]0FF4}{[1][0][]source}{}{}{}}
\makeatother
\crefname{originalreference}{original reference}{original references}
\Crefname{originalreference}{Original reference}{Original references}
\makeatletter
\newlabel{lemma-functor-and-cycles}{{0FG5}{source}{}{sourcepointer-3}{}}
\newlabel{lemma-functor-and-cycles@cref}{{[originalreference][0][]0FG5}{[1][0][]source}{}{}{}}
\makeatother
\crefname{originalreference}{original reference}{original references}
\Crefname{originalreference}{Original reference}{Original references}
\makeatletter
\newlabel{lemma-manin}{{0FGN}{source}{}{sourcepointer-4}{}}
\newlabel{lemma-manin@cref}{{[originalreference][0][]0FGN}{[1][0][]source}{}{}{}}
\makeatother
\title{The projective space bundle formula for motives}
\author{The Stacks Project Authors}
\date{stacks-project:weil.tex:local-snapshot-20261012; GNU FDL 1.2 or later}
\begin{document}
\maketitle
\noindent\textit{Editorial note.} All delivered mathematical text is the author's own, copied from the pinned Stacks Project chapter file weil.tex; no statement or proof step was rewritten, reordered or removed.\par\medskip
\noindent\textit{Editorial note.} The counted proof is the author's complete argument: Prefilter hit re-read per the 2026-10-14 12:4x THIN-criterion patch (pb>=800 / pw>=100): the proof body was read in full at source and its main line is complete; the flagged phrase is incidental and is quoted verbatim below. Verbatim, end of proof line 1415/1416: 'Some details omitted.'  -  the details are the routine identification of the pushforwards in the standard formula; the reduction and the identification of the map are given in the span.\par\medskip
\noindent\textit{Editorial note.} The label lemma-projective-space-bundle-formula names the counted statement, and the proof environment immediately following it in the source is delivered verbatim as its proof.\par\medskip
\noindent\textit{Editorial note.} The counted statement is carried by the author's own lemma environment in the chapter source.\par\medskip
\noindent\textit{Editorial note.} Prefilter hit re-read per the 2026-10-14 12:4x THIN-criterion patch: proof body (proof\_lines 1402-1416, pb=776, pw=109) read in full; main line complete.\par\medskip
\noindent\textit{Editorial note.} Editorial note (verbatim): 'Some details omitted.' at the end of the proof, following the explicit identification of the pushforwards along the Chern class correspondences.\par\medskip
\noindent\textit{Editorial note.} Pre-open gate: python operations/agent-authoring/tools/precard\_gate.py --source weil.tex --label lemma-projective-space-bundle-formula --start 1390 --end 1416 --local-id 28164 returned RESULT=OK. Pre-assembly tag self-check: canonical tag resolved and not present in reference/public-items.jsonl.\par\medskip
\noindent\textit{Editorial note.} The statement and proof are copied from the pinned chapter file; the cited inputs are kept as external pointers with their canonical Stacks tags.\par\medskip
\noindent\textit{Allowed standard prerequisites (editorial).}\begin{itemize}
\item Manin's criterion for isomorphisms of motives. Same chapter, Manin's criterion, reducing the statement to an isomorphism on Chow groups after tensoring with the motive of an arbitrary smooth projective scheme. Reduces the isomorphism of motives to an isomorphism of Chow groups after an arbitrary base change.
\item The projective space bundle formula for Chow groups and the description of pushforwards along graphs. Chow Homology, the projective space bundle formula, and Same chapter, the description of pushforward of cycles along the correspondence of a graph as pullback of cycles. Identifies the map on Chow groups with the standard formula for the projectivization of a vector bundle.
\end{itemize}\medskip
\ifdefined\hypertarget\hypertarget{sourcepointer-1}{}\fi
\noindent\textit{Source reference.} Original lemma 02TX: sources/stacks-chow/tags.txt, source lines 3712--3712. Prerequisite pointer invoked at proof line(s) 1410; it names a standard fact of a sibling chapter that the counted argument uses, and every step of the counted proof built on it is retained verbatim in the delivered body.\par\medskip
\ifdefined\hypertarget\hypertarget{sourcepointer-2}{}\fi
\noindent\textit{Source reference.} Original lemma 0FF4: sources/stacks-chow/tags.txt, source lines 19274--19274. Prerequisite pointer invoked at proof line(s) 1413; it names a standard fact of a sibling chapter that the counted argument uses, and every step of the counted proof built on it is retained verbatim in the delivered body.\par\medskip
\ifdefined\hypertarget\hypertarget{sourcepointer-3}{}\fi
\noindent\textit{Source reference.} Original reference 0FG5: sources/stacks-weil/pinned-original-source.tex, source lines 661--661. Prerequisite pointer invoked at proof line(s) 1412; it names a standard fact of the same chapter that the counted argument uses, and every step of the counted proof built on it is retained verbatim in the delivered body.\par\medskip
\ifdefined\hypertarget\hypertarget{sourcepointer-4}{}\fi
\noindent\textit{Source reference.} Original reference 0FGN: sources/stacks-weil/pinned-original-source.tex, source lines 1327--1327. Prerequisite pointer invoked at proof line(s) 1403; it names a standard fact of the same chapter that the counted argument uses, and every step of the counted proof built on it is retained verbatim in the delivered body.\par\medskip
\begin{lemma}[Projective space bundle formula]
\label{lemma-projective-space-bundle-formula}
In the situation above, the map
$$
\sum\nolimits_{i = 0, \ldots, r - 1} c_i :
\bigoplus\nolimits_{i = 0, \ldots, r - 1} h(X)(-i)
\longrightarrow
h(P)
$$
is an isomorphism in the category of motives.
\end{lemma}

\begin{proof}
By Lemma \ref{lemma-manin} it suffices to show that
our map defines an isomorphism on Chow groups of motives
after taking the product with any smooth projective scheme $Z$.
Observe that $P \times Z \to X \times Z$ is the projective
bundle associated to the pullback of $\mathcal{E}$ to $X \times Z$.
Hence the statement on Chow groups is true by the projective space bundle
formula given in
Chow Homology, Lemma \ref{chow-lemma-chow-ring-projective-bundle}.
Namely, pushforward of cycles along $[\Gamma_p]$ is given by pullback
of cycles by $p$ according to Lemma \ref{lemma-functor-and-cycles} and
Chow Homology, Lemma \ref{chow-lemma-lci-gysin-flat}. Hence pushforward
along $c_i$ sends $\alpha$ to $c_1(\mathcal{O}_P(1))^i \cap p^*\alpha$.
Some details omitted.
\end{proof}

\end{document}
